%%%%%%%%%%%%%%%%%%%%%%%%%%%%%%%%%%%%%%%%%%%%%%%%%%%%%%%%%%%%%%%%%%%%%%%%%%%%%%%%%%%%%%%%%%%%%%%%%%%%%%%
%%%%%%%%%%%%%% Template de Artigo Adaptado para Trabalho de Conclusão de Curso - SI Contagem - PUCMINAS %%%%%%%%%%%%%%%%%%%%%%%%
%% codificação UTF-8 - Abntex - Latex -  							     %%
%% Autor da primeira versão:    Fábio Leandro Rodrigues Cordeiro  (fabioleandro@pucminas.br)                            %% 
%% Co-autor da primeira versão: Prof. João Paulo Domingos Silva, Harison da Silva e Anderson Carvalho                   %%
%% Revisores normas NBR (Padrão PUC Minas) da primeira versão: Helenice Rego Cunha e Prof. Theldo Cruz                  %%
%% Versão: 1.1     18 de dezembro 2015                     	                                     %%
%%%%%%%%%%%%%%%%%%%%%%%%%%%%%%%%%%%%%%%%%%%%%%%%%%%%%%%%%%%%%%%%%%%%%%%%%%%%%%%%%%%%%%%%%%%%%%%%%%%%%%%


\documentclass[a4paper,12pt]{article}
\usepackage{times}
\usepackage{abakos}  %pacote com padrão da Abakos baseado no padrão da PUC
%%%%%%%%%%%%%%%%%%%%%%%%%%%
%Capa da revista
%%%%%%%%%%%%%%%%%%%%%%%%%%

%\setcounter{page}{80} %iniciar contador de pagina de valor especificado
\newcommand{\monog}{Os Tipos Lógicos de Pesquisa: Uma Síntese Baseada na Obra de Gilson Volpato}
%\newcommand{\monogES}{Article template Institute of Mathematical Sciences and Informatics}
\newcommand{\tipo}{Artigo}  % Especificar a seção tipo do trabalho: Artigo, Resumo, Tese, Dociê etc
%\newcommand{\origem}{Brasil}
%\newcommand{\editorial}{Belo Horizonte, p. 01-11, nov. 2024}  % p. xx-xx – páginas inicial-final do artigo
\newcommand{\editorial}{}  
\newcommand{\lcc}{\scriptsize{Licença Creative Commons Attribution-NonCommercial-NoDerivs 3.0 Unported}}

%%%%%%%%%%%%%%%%%INFORMAÇÕES SOBRE AUTOR PRINCIPAL %%%%%%%%%%%%%%%%%%%%%%%%%%%%%%%
\newcommand{\AutorA}{Samuel Teodoro Albuquerque Fragoso}
\newcommand{\funcaoA}{}
\newcommand{\emailA}{\href{mailto:samuel.fragoso@sga.pucminas.br}{samuel.fragoso@sga.pucminas.br}}
\newcommand{\cursA}{Aluno(a) do Programa de Graduação em Ciência da Computação}

%\newcommand{\AutorB}{Nome completo do(a) orientador(a)}
%\newcommand{\funcaoB}{}
%\newcommand{\emailB}{\href{mailto:xxx@sga.pucminas.br}{xxx@pucminas.br}}
%\newcommand{\cursB}{Professor(a) do Programa de Graduação em Sistemas de Informação}
% 
% Definir macros para o nome da Instituição, da Faculdade, etc.
\newcommand{\univ}{Pontifícia Universidade Católica de Minas Gerais}

\newcommand{\keyword}[1]{\textsf{#1}}

%Para as URLS não ultrapassarem a margem
\usepackage{url}
\makeatletter
\g@addto@macro{\UrlBreaks}{\UrlOrds}
\makeatother

\usepackage{adjustbox} % reduzir tamanho figuras

\usetikzlibrary{arrows,shapes,positioning,shadows,trees}

\tikzset{
  basic/.style  = {draw, text width=3cm, drop shadow, font=\sffamily, rectangle},
  root/.style   = {basic, rounded corners=2pt, thin, align=center,
                   fill=blue!30},
  level 2/.style = {basic, rounded corners=6pt, thin,align=center, fill=green!30,
                   text width=8em},
  level 3/.style = {basic, thin, align=left, fill=pink!60, text width=6.5em}
}

\begin{document}
% %%%%%%%%%%%%%%%%%%%%%%%%%%%%%%%%%%
% %% Pagina de titulo
% %%%%%%%%%%%%%%%%%%%%%%%%%%%%%%%%%%

\begin{center}
\includegraphics[scale=0.2]{figuras/brasao.jpg} \\
PONTIFÍCIA UNIVERSIDADE CATÓLICA DE MINAS GERAIS \\
Instituto de Ciências Exatas e de Informática

% \vspace{1.0cm}

\end{center}

 \vspace{0cm} {
 \singlespacing \Large{\monog \symbolfootnote[1]{Artigo apresentado ao Instituto de Ciências Exatas e Informática da Pontifícia Universidade Católica de \linebreak Minas Gerais, campus Coração Eucarístico.} \\ }
 % \normalsize{\monogES}
 }

\vspace{1.0cm}

\begin{flushright}
\singlespacing 
\normalsize{\AutorA \footnote{\funcaoA \cursA \,-- \emailA . }} \\
%\normalsize{\AutorC \footnote{\funcaoC \cursC \,-- \emailC . }} \\
%\normalsize{\AutorD \footnote{\funcaD \cursD \,-- \emailD . }} \\
\end{flushright}
\thispagestyle{empty}

\vspace{1.0cm}

\begin{abstract}
\noindent
O vídeo de Gilson Volpato elabora que existem 3 tipos de pesquisas, a primeira que não precisa de hipóteses para ter suas perguntas respondidas e as outras duas que precisam da elaboração de hipóteses para guiar o pesquisador ao resultado. Os dois tipos que precisam de hipóteses são divididos em mais dois tipos, que são associação e interferência, as duas testam a relação entre duas ou mais variáveis. No caso da associação é definida quando há associação mas não efito e a interferência quando é encontrada uma relação entre causa e efeito. 
\\\textbf{\keyword{Palavras-chave:}} Metodologia Científica; Tipos de Pesquisa; Hipótese; Causalidade.
\end{abstract}

%%%%%%%%%%%%%%%%%%%%%%%%%%%%%%%%%%%%%%%%%%%%%%%%%%%%%%%%%
\selectlanguage{english}
\begin{abstract}
\noindent
Gilson Volpato's video explains that there are 3 types of research: the first one does not require hypotheses to have its questions answered, and the other two require the elaboration of hypotheses to guide the researcher to the result. The two types that need hypotheses are divided into two further types, which are association and inference; both test the relationship between two or more variables. Association is defined when there is a relationship but no effect, and inference is when a cause-and-effect relationship is found.
\\\textbf{\keyword{Keywords:}} Scientific Methodology; Research Types; Hypothesis; Causality.
\end{abstract}

\selectlanguage{brazilian}
 \onehalfspace  % espaçamento 1.5 entre linhas
 \setlength{\parindent}{1.25cm}

%%%%%%%%%%%%%%%%%%%%%%%%%%%%%%%%%%%%%%%%%%%%%%%%%
%% INICIO DO TEXTO
%%%%%%%%%%%%%%%%%%%%%%%%%%%%%%%%%%%%%%%%%%%%%%%%%

%%%%%%%%%%%%%%%%%%%%%%%%%%%%%%%%%%%%%%%%%%%%%%%%%%%%%%%%%%%%%%%%%%%%%%%%%%%%%%%%%%%%%%%%%%%%%%%%%%%%%%%%%%%%%%%%%%%%%%%%%%%%%%%%
%%%%%%%%%%%%%% Template de Artigo Adaptado para Trabalho de Conclusão de Curso - SI Contagem - PUCMINAS                       %%
%% codificação UTF-8 - Abntex - Latex -  							                                                          %%
%% Autor da primeira versão:    Fábio Leandro Rodrigues Cordeiro                                                              %% 
%% Co-autores da primeira versão: Prof. João Paulo Domingos Silva, Harison da Silva e Anderson Carvalho		                  %%
%% Revisores normas NBR (Padrão PUC Minas) da primeira versão: Helenice Rego Cunha e Prof. Theldo Cruz                        %%
%% Versão: 1.1     18 de dezembro 2015                                                                                        %%
%%%%%%%%%%%%%%%%%%%%%%%%%%%%%%%%%%%%%%%%%%%%%%%%%%%%%%%%%%%%%%%%%%%%%%%%%%%%%%%%%%%%%%%%%%%%%%%%%%%%%%%%%%%%%%%%%%%%%%%%%%%%%%%%
\section{\esp Introdução (1 página)} 

O presente trabalho tem como objetivo apresentar uma síntese das classificações de pesquisa científica abordadas por Gilson Volpato\cite{volpato2012}. Será feita a distinção entre os tipos lógicos de pesquisa, com foco na diferenciação entre pesquisas que necessitam da criação de hipóteses - como no caso da associação e interferência - e das que não possuem essa exigência para atingir o resultado.


\section{\esp Referencial Teórico (1 página)}
Deixando de lado as especifidades e qualificações de uma pesquisa e entrando em sua raíz lógica podemos defini-la em até três tipos. Entender esses três tipos é fundamental para todo o processo de fazer uma pesquisa e escrever um artigo. A pesquisa se inicia com uma pergunta sobre um fato observado, essa resposta pode ser, em dois dos três tipos, respondida com o auxílio de hipóteses científicas. De acordo com a Encyclopaedia Britannica (2024)\cite{britannica2024}, uma hipótese científica é uma explicação inicial para um fato observado, que ainda precisa ser testada. No primeiro caso a pergunta não precisa necessariamente de uma hipótese para ser respondida como no exemplo da pergunta "Que rotas de migração são utilizadas por uma população de baleias jubarte?". Aqui o objetivo é rastrear e descrever o deslocamento, é uma questão de "o quê" e "onde", não "por quê" e por isso não se faz necessário o uso de uma hipótese. A coleta de dados é direcionada apenas pela pergunta e não por uma hipótese que nesse caso não passaria de um capricho acadêmico.  A hipótese se faz necessária quando a pergunta da pesquisa busca explicar ou comparar. No exemplo da pergunta "Por que a luz não acende?" propõe-se a hipótese de que "A lâmpada queimou", depois de testá-la em um soquede e descobrir que estava errada, formula-se uma nova hipótese: "interruptor quebrado". Com a realização de novos testes, a hipótese é confirmada. Nesse exemplo a coleta de dados foi direcionada pela hipótese e a elaboração dela foi crucial para a descoberta do resultado, mostrando-se uma ferramenta fundamental para chegar na resposta do problema. É importante que as primeiras hipóteses sejam as mais simples e se forem refutadas inicia-se os testes de hipótese mais complexas. As pesquisas com hipóteses examinam a correlação ou causalidade entre variáveis e são dividas em dois tipos: as de associação e interferência. A associação é uma relação onde duas variáveis mudam de forma previsível, mas sem uma relação de causa e efeito direta. Como é mostrado no artigo "New evidence for the Theory of the Stork"\cite{hofer2004}, de Thomas Höfer, onde cita dados de Berlim que mostram uma correlação entre crescimento do número de cegonhas (variável A) na cidade com partos fora de hospitais (variável B) mas não há qualquer relação entre hospitais e população de cegonhas, mas pode haver uma variável C oculta influenciando ambas, caracterizando uma relação sem causa e efeito entre elas. A interferência ocorre quando há uma relação de causa, mecanismo e efeito onde as variáveis estão relacionadas e existe uma relação de causa e efeito sobre elas. Portanto, os três tipos lógicos de pesquisa são: descrição, associação e associação com interferência.




\section{\esp Conclusão}


Este trabalho apresentou a distinção fundamental entre os três tipos lógicos de pesquisa: descritiva, de associação e de interferência. Foi demonstrado que a principal diferença entre eles reside na necessidade e no uso de hipóteses para guiar a investigação. A compreensão dessa estrutura é essencial para a elaboração de projetos de pesquisa claros e metodologicamente corretos, garantindo que as perguntas, os métodos e as conclusões de um estudo científico estejam alinhados de forma lógica e coerente.




%%%%%%%%%%%%%%%%%%%%%%%%%%%%%%%%%%%
%% FIM DO TEXTO
%%%%%%%%%%%%%%%%%%%%%%%%%%%%%%%%%%%

% \selectlanguage{brazil}
%%%%%%%%%%%%%%%%%%%%%%%%%%%%%%%%%%%
%% Inicio bibliografia
%%%%%%%%%%%%%%%%%%%%%%%%%%%%%%%%%%%

 \newpage
\singlespace{
\renewcommand\refname{REFERÊNCIAS}
\bibliography{bibliografia}
}

%Inclusão do arquivo abntex2-alf.bst como solução para adequação à ABNT NBR 10520:2023 quanto às citações, que não são mais em caixa-alta
\bibliographystyle{abntex2-alf.bst}

\end{document}